\documentclass[12pt]{article}
\usepackage{amsmath,amssymb,amsfonts,mathrsfs}
\usepackage[margin=1in]{geometry}

\begin{document}

\section*{Chapter 2, Problem 11}

\subsection*{Problem Statement}
The Lagrangian density $\mathscr{L}$ which generates a given set of Euler--Lagrange equations is not unique. Prove this result by showing that adding a divergence to $\mathscr{L}$ does not alter the Euler--Lagrange equations. That is, let
\[
  \mathscr{L}' = \mathscr{L} + \sum_{k=1}^{n} \frac{\partial \Lambda_k}{\partial x_k},
\]
where
\[
  \mathscr{L}' = \mathscr{L}'\!\left(x_k, \psi_j, \frac{\partial\psi_j}{\partial x_k}\right), \quad \Lambda_k = \Lambda_k\!\left(x_k, \psi_j\right),
\]
$j = 1,\ldots,m$ indexes the dependent field variables, and $k = 1,\ldots,n$ indexes the independent variables.

Now prove that
\[
  \frac{\partial\mathscr{L}'}{\partial\psi_j} - \sum_k \frac{\partial}{\partial x_k}\frac{\partial\mathscr{L}'}{\partial(\partial\psi_j/\partial x_k)} = \frac{\partial\mathscr{L}}{\partial\psi_j} - \sum_k \frac{\partial}{\partial x_k}\frac{\partial\mathscr{L}}{\partial(\partial\psi_j/\partial x_k)},
\]
where we define
\[
  \frac{\partial\Lambda_k}{\partial x_k} = \frac{\partial\Lambda_k}{\partial x_k}\bigg|_{\text{explicit}} + \sum_j \frac{\partial\Lambda_k}{\partial\psi_j}\frac{\partial\psi_j}{\partial x_k}.
\]

\subsection*{Solution}

We need to show that the divergence term $\sum_k \partial\Lambda_k/\partial x_k$ does not contribute to the Euler--Lagrange equations.

Since $\Lambda_k = \Lambda_k(x_k, \psi_j)$ (depending only on the coordinates and fields, not on field derivatives), the total divergence is:
\[
  \Phi \equiv \sum_{k=1}^{n}\frac{\partial\Lambda_k}{\partial x_k}
  = \sum_k\left[\frac{\partial\Lambda_k}{\partial x_k}\bigg|_{\text{expl}} + \sum_j \frac{\partial\Lambda_k}{\partial\psi_j}\frac{\partial\psi_j}{\partial x_k}\right].
\]

We compute the Euler--Lagrange expression for $\Phi$ with respect to $\psi_j$:
\[
  E_j(\Phi) \equiv \frac{\partial\Phi}{\partial\psi_j} - \sum_k \frac{\partial}{\partial x_k}\frac{\partial\Phi}{\partial(\partial\psi_j/\partial x_k)}.
\]

\medskip
\noindent\textbf{Step 1:} Compute $\dfrac{\partial\Phi}{\partial(\partial\psi_j/\partial x_k)}$.

Since $\Phi = \sum_k\left[\frac{\partial\Lambda_k}{\partial x_k}\big|_{\text{expl}} + \sum_\ell \frac{\partial\Lambda_k}{\partial\psi_\ell}\frac{\partial\psi_\ell}{\partial x_k}\right]$, and $\Lambda_k$ does not depend on field derivatives, the only dependence on $\partial\psi_j/\partial x_k$ comes from the second sum. Thus (note the index $k$ in $\Lambda_k$ matches with $x_k$, so for a fixed $k$):
\[
  \frac{\partial\Phi}{\partial(\partial\psi_j/\partial x_k)} = \frac{\partial\Lambda_k}{\partial\psi_j}.
\]

\medskip
\noindent\textbf{Step 2:} Compute $\dfrac{\partial\Phi}{\partial\psi_j}$.

\[
  \frac{\partial\Phi}{\partial\psi_j} = \sum_k \left[\frac{\partial^2\Lambda_k}{\partial\psi_j\partial x_k}\bigg|_{\text{expl}} + \sum_\ell \frac{\partial^2\Lambda_k}{\partial\psi_j\partial\psi_\ell}\frac{\partial\psi_\ell}{\partial x_k}\right].
\]

\medskip
\noindent\textbf{Step 3:} Compute $\displaystyle\sum_k \frac{\partial}{\partial x_k}\frac{\partial\Phi}{\partial(\partial\psi_j/\partial x_k)} = \sum_k \frac{\partial}{\partial x_k}\frac{\partial\Lambda_k}{\partial\psi_j}$.

Expanding the total derivative:
\[
  \sum_k \frac{\partial}{\partial x_k}\frac{\partial\Lambda_k}{\partial\psi_j}
  = \sum_k\left[\frac{\partial^2\Lambda_k}{\partial x_k\partial\psi_j}\bigg|_{\text{expl}} + \sum_\ell\frac{\partial^2\Lambda_k}{\partial\psi_\ell\partial\psi_j}\frac{\partial\psi_\ell}{\partial x_k}\right].
\]

\medskip
\noindent\textbf{Step 4:} Subtract:
\[
  E_j(\Phi) = \frac{\partial\Phi}{\partial\psi_j} - \sum_k\frac{\partial}{\partial x_k}\frac{\partial\Lambda_k}{\partial\psi_j}.
\]

By comparing Steps 2 and 3, we see they are identical (assuming sufficient smoothness so that mixed partial derivatives commute):
\[
  E_j(\Phi) = \sum_k\left[\frac{\partial^2\Lambda_k}{\partial\psi_j\partial x_k} + \sum_\ell\frac{\partial^2\Lambda_k}{\partial\psi_j\partial\psi_\ell}\frac{\partial\psi_\ell}{\partial x_k}\right]
  - \sum_k\left[\frac{\partial^2\Lambda_k}{\partial x_k\partial\psi_j} + \sum_\ell\frac{\partial^2\Lambda_k}{\partial\psi_\ell\partial\psi_j}\frac{\partial\psi_\ell}{\partial x_k}\right] = 0.
\]

Therefore:
\[
  \frac{\partial\mathscr{L}'}{\partial\psi_j} - \sum_k\frac{\partial}{\partial x_k}\frac{\partial\mathscr{L}'}{\partial(\partial\psi_j/\partial x_k)}
  = \frac{\partial\mathscr{L}}{\partial\psi_j} - \sum_k\frac{\partial}{\partial x_k}\frac{\partial\mathscr{L}}{\partial(\partial\psi_j/\partial x_k)} + E_j(\Phi)
  = \frac{\partial\mathscr{L}}{\partial\psi_j} - \sum_k\frac{\partial}{\partial x_k}\frac{\partial\mathscr{L}}{\partial(\partial\psi_j/\partial x_k)}.
\]

This proves that adding a divergence to the Lagrangian density does not change the Euler--Lagrange equations.

\medskip
\noindent\textbf{Physical interpretation:} This result can also be understood from the action principle. Adding $\sum_k \partial\Lambda_k/\partial x_k$ to $\mathscr{L}$ adds to the action a term $\int \sum_k \partial\Lambda_k/\partial x_k\,d^n x$, which by the divergence theorem equals a surface integral $\oint \Lambda_k\,dS_k$ over the boundary. Since the field values are fixed on the boundary (as boundary conditions), this surface integral is a constant under variations, and hence does not affect the extremum condition.

\end{document}
